%%%%%%%%%%%%%%%%%%%%%%%%%%%%%%%%%%%%%%%%%%%%%%%%%%%%%%%%%%%%%%%%%%%%%%%%%%%%%%%%%%%%%%
%%%%%%%%%%%%%%%%%%%%  INGRESAR LAS PREGUNTAS DE LA  EVALUACIÓN %%%%%%%%%%%%%%%%%%%%%%%%
%%%%%%%%%%%%%%%%%%%%%%%%%%%%%%%%%%%%%%%%%%%%%%%%%%%%%%%%%%%%%%%%%%%%%%%%%%%%%%%%%%%%%%%

\paragraph{I} Multiple option section. 
Read the following questions and identify the option that best answers each, then write the letter on the line.

\begin{preguntas}

\pregunta[4] \rule{1cm}{0.1mm} Identify the quotient rule for derivatives. 

\begin{enumerate*}[label=(\alph*)]
    \item \(\left(\dfrac{u}{v}\right)' = \dfrac{v'u - u'v}{u^2}\)\hspace{2em}
    \item \(\left(\dfrac{u}{v}\right)' = \dfrac{u'v - uv'}{v^2}\) \hspace{2em}
    \item \(\left(\dfrac{u}{v}\right)' = \dfrac{uv' - u'v}{v^2}\) \hspace{2em} 
    \item \(\left(\dfrac{u}{v}\right)' = \dfrac{u'v - v'u}{u^2}\). 
\end{enumerate*}

\pregunta[4] \rule{1cm}{0.1mm} Identify the product rule for derivatives 

\begin{enumerate*}[label=(\alph*)]
    \item \( (uv)' = u' + v' \)\hspace{2em}
    \item \( (uv)' = u'v + uv' \) \hspace{2em}
    \item \( (uv)' = uv' + uv' \)\hspace{2em}
    \item \( (uv)' = u'v + u'v \)
\end{enumerate*}

\pregunta[4] \rule{1cm}{0.1mm} Identify the chain rule for derivatives 

\begin{enumerate*}[label=(\alph*)]
    \item \( \left[u(v(x))\right]' = v'(u(x)) \) 
    \item \( \left[u(v(x))\right]' = v'(u(x))u'(x) \)
    \item \( \left[u(v(x))\right]' = u'(v(x))v'(x) \)
    \item \( \left[u(v(x))\right]' = u'(v(x)) \)
\end{enumerate*}

\pregunta[4] \rule{1cm}{0.1mm} Which of the following derivatives is correct
%Select the correct type of function of $f(x):=x^n$, where $n\in\mathbb{Z}$ ($x$ belongs to the set of integer numbers

\begin{enumerate*}[label=(\alph*)]
    \item \(\dv{x}n^{x} = n^{x}\ln(n)\)\hspace{2em}
    \item \(\dv{x}n^{x} = \frac{1}{x}n^{x}\)\hspace{2em}
    \item \(\dv{x}n^{x} = xn^{x}\) \hspace{2em}
    \item \(\dv{x}n^{x} = ne^{x-1}\)
\end{enumerate*}

\pregunta[4] \rule{1cm}{0.1mm} Which of the following derivatives is correct
%Select the option of the function that map all the domain into the positive real numbers.

\begin{enumerate*}[label=(\alph*)]
    \item \(\dv{x}\sin(x) = \cos(x)\) \hspace{2em}
    \item \(\dv{x}\sin(x) = -\sin(x)\)\hspace{2em}
    \item \(\dv{x}\sin(x) = \sin(x)\)\hspace{2em}
    \item \(\dv{x}\sin(x) = -\cos(x)\)
\end{enumerate*}

\pregunta[4] \rule{1cm}{0.1mm} Which of the following derivatives is correct

\begin{enumerate*}[label=(\alph*)]
    \item \(\dv{x}\cos(x) = \cos(x)\) \hspace{2em}
    \item \(\dv{x}\cos(x) = -\sin(x)\)\hspace{2em}
    \item \(\dv{x}\cos(x) = \sin(x)\)\hspace{2em}
    \item \(\dv{x}\cos(x) = -\cos(x)\)
\end{enumerate*}

\pregunta[4] \rule{1cm}{0.1mm} Which of the following derivatives is correct
%Select the option of the function that map all the domain into the positive real numbers.

\begin{enumerate*}[label=(\alph*)]
    \item \(\dv{x}\log_n(x) = \dfrac{1}{\ln(x)}\)\hspace{2em} 
    \item \(\dv{x}\log_n(x) = \dfrac{1}{x\ln(n)}\)\hspace{2em}
    \item \(\dv{x}\log_n(x) = \dfrac{\ln(x)}{x}\)\hspace{2em}
    \item \(\dv{x}\log_n(x) = \dfrac{1}{x\ln(x)}\)
\end{enumerate*}

\pregunta[4] \rule{1cm}{0.1mm} Which of the following options is the implicit derivative of \(y^2 = x\) 
%Select the option of the function that map all the domain into the positive real numbers.

\begin{enumerate*}[label=(\alph*)]
    \item \(\dv{x}{y}=0\)\hspace{2em} 
    \item \(\dv{x}{y}=2y-1\)\hspace{2em}
    \item \(\dv{x}{y}=1\)\hspace{2em}
    \item \(\dv{x}{y}=1-2y\)
\end{enumerate*}



\end{preguntas}

\newpage

\paragraph{II} Implicit derivatives.
Use the rules of derivatives that can be usefull to compute the implicit derivatives of the following functions. 
Compute \(\dv{y}{x}\) and \(\dv{y}{x}\).

\begin{preguntas}

\pregunta[10]{Compute the implicit derivative of the function \(f(x) = e^{x + y}= x^2y^3 \).}

\begin{center}
    \framebox(\textwidth,250){} 
\end{center}

\pregunta[10]{Compute the implicit derivative of the function \(g(x) = x^3+y^3=6xy \).}

\begin{center}
    \framebox(\textwidth,250){} 
\end{center}

\end{preguntas}

\newpage

\paragraph{III} Critical points.
In this section your task is to find the critical points of implicit derivatives considering only \(\dv{y}{x}=0\).

\begin{preguntas}

\pregunta[13]{Find the critical points of \(x^{2/3}+y^{2/3}=4^{2/3}\)}

\begin{center}
    \framebox(\textwidth,250){} 
\end{center}

\pregunta[13]{Find the crictical points of \(\dfrac{(x-3)^2}{9}-\dfrac{(y-8)^{2}}{4}=1\)}

\begin{center}
    \framebox(\textwidth,250){} 
\end{center}

\newpage 

\paragraph{IV} High order derivatives.
Compute the second order derivative of the following functions and reduce them as muach as possible.

\pregunta[11]{\( f(x) = e^{x+3}\)}

\begin{center}
    \framebox(\textwidth,250){} 
\end{center}

\pregunta[11]{\( h(x) = e^{3x}\sin(\dfrac{\pi}{3}x ) \)}

\begin{center}
    \framebox(\textwidth,250){} 
\end{center}



\end{preguntas}

